\documentclass[11pt]{article}
\usepackage[T1]{fontenc}
\usepackage[french]{babel}
\usepackage[margin=2.5cm, headheight=14pt]{geometry}
\usepackage{amsmath, amssymb, amsthm}
\usepackage{enumitem}
\usepackage{fancyhdr}

\newcommand{\cours}{SIGLE1234 -- Nom du cours}
\newcommand{\etudiant}{{{AUTHOR}}}

\pagestyle{fancy}
\fancyhf{}
\fancyhead[L]{\cours}
\fancyhead[R]{\etudiant}
\fancyfoot[C]{\thepage}

\theoremstyle{definition}
\newtheorem{exercice}{Exercice}
\newenvironment{solution}{\begin{proof}[Solution]}{\end{proof}}

\begin{document}

\begin{center}
  {\LARGE\bfseries {{TITLE}} \par}
  \medskip
  \etudiant \quad -- \quad Remis le \today
\end{center}

\begin{exercice}
Montrer que pour tout entier $n \ge 1$,
\[
  \sum_{k=1}^{n} k = \frac{n(n+1)}{2}.
\]
\end{exercice}

\begin{solution}
Par récurrence sur $n$. Pour $n = 1$, les deux membres valent~$1$.
Supposons l'égalité vraie au rang $n$. Alors
\[
  \sum_{k=1}^{n+1} k = \frac{n(n+1)}{2} + (n+1) = \frac{(n+1)(n+2)}{2},
\]
ce qui est l'égalité au rang $n+1$.
\end{solution}

\begin{exercice}
Soit $f(x) = x^3 - 3x + 1$.
\begin{enumerate}[label=(\alph*)]
  \item Calculer $f'(x)$.
  \item Trouver les extremums locaux de $f$.
\end{enumerate}
\end{exercice}

\begin{solution}
\begin{enumerate}[label=(\alph*)]
  \item $f'(x) = 3x^2 - 3 = 3(x - 1)(x + 1)$.
  \item $f'$ s'annule en $x = -1$ et en $x = 1$. Comme $f''(x) = 6x$, $f$
        admet un maximum local $f(-1) = 3$ et un minimum local $f(1) = -1$.
\end{enumerate}
\end{solution}

\end{document}
